\documentclass[12pt]{article}

\usepackage[spanish]{babel}
\usepackage[utf8]{inputenc}
\usepackage[T1]{fontenc}
\usepackage{geometry}
\usepackage{setspace}
\usepackage{array}
\usepackage{titlesec}
\usepackage{tabularx}

\geometry{letterpaper, margin=2.5cm}
\setstretch{1.3}

\begin{document}

\begin{flushright}
San Luis Potosí, S.L.P. a 20 de Octubre de 2026
\end{flushright}

\vspace{1cm}

\textbf{ASUNTO:} Solicitud de horas de clase para reuniones del AWS Student Group.

\vspace{1cm}

\textbf{COORDINACIÓN DE LA CARRERA DE} \\
\textbf{INGENIERÍA EN SISTEMAS COMPUTACIONALES} \\
P R E S E N T E.

\vspace{1cm}

Por medio de la presente, solicitamos atentamente su autorización para llevar a cabo una serie de talleres prácticos en el marco del AWS Student Group del Instituto Tecnológico de Matehuala, durante el periodo comprendido entre el 25 de octubre y el 7 de diciembre de 2026.

El objetivo de estas actividades es fortalecer los conocimientos teóricos y prácticos de los estudiantes en materia de computación en la nube, administración de sistemas Linux, infraestructura como servicio, despliegue de aplicaciones web y desarrollo de soluciones basadas en servicios de Amazon Web Services (AWS).

Para ello, proponemos un horario de actividades, las cuales según el Departamento de Sistemas pueden realizarse en el horario que nos puedan asignar de 7:00 AM a 1:00 PM, según la disponibilidad de los profesores, divididas en sesiones consecutivas de una hora tras otra en 2 grupos (de acuerdo con la cantidad de alumnos que se presenten), en las fechas establecidas en el calendario de actividades anexo a la presente solicitud. Las sesiones tendrán un enfoque práctico, con ejercicios guiados y el desarrollo progresivo de proyectos tecnológicos.

Asimismo, nos comprometemos a hacer un uso responsable de las instalaciones y recursos institucionales, así como a respetar los lineamientos académicos y administrativos correspondientes.

\vspace{0.5cm}
\textbf{Nota sobre el Calendario:} Se informa que la Sesión 1 (Introducción a AWS) ya fue impartida con anterioridad. El siguiente calendario es de carácter \textbf{tentativo} y está sujeto a las modificaciones que el Departamento de Sistemas y la Coordinación de la Carrera consideren pertinentes, con el fin de adaptarse a la disponibilidad de aulas y horarios de los docentes. Se mantiene como fecha fija de clausura el día \textbf{7 de diciembre de 2026}.
\vspace{0.3cm}
\begin{center}
\begin{tabularx}{\textwidth}{|c|c|X|}
\hline
\textbf{Semana} & \textbf{Fechas} & \textbf{Tema Principal} \\
\hline
1 & 26 y 28 de Octubre & Sesión 2: Despliegue Web (S3 y CloudFront) \\
\hline
2 & 2 y 4 de Noviembre & Sesión 3: Servidores Linux (EC2 y SSH) \\
\hline
3 & 11 y 18 de Noviembre & Sesión 4: Servidor de Minecraft (Puertos y Almacenamiento) \\
\hline
4 & 23 y 25 de Noviembre & Sesión 5: APIs sin Servidor (Lambda y DynamoDB) \\
\hline
5 & 30 de Nov. y 2 de Dic. & Práctica libre y repaso (Por definir con Sistemas) o charla sobre software libre y AWS \\
\hline
6 & 7 de Diciembre & Clausura del programa y entrega de reconocimientos (Pequeño Convivio) \\
\hline
\end{tabularx}
\end{center}
\vspace{0.5cm}
\textbf{Nota sobre el Calendario:} Se informa que la Sesión 1 (Introducción a AWS) ya fue impartida con anterioridad. El siguiente calendario es de carácter \textbf{tentativo} y está sujeto a las modificaciones que el Departamento de Sistemas y la Coordinación de la Carrera consideren pertinentes, con el fin de adaptarse a la disponibilidad de aulas y horarios de los docentes. Se mantiene como fecha fija de clausura el día \textbf{7 de diciembre de 2026}.
\textbf{Requisito indispensable para los alumnos:} Para garantizar el óptimo aprovechamiento de los talleres y agilizar el proceso de las prácticas en clase, es \textbf{fundamental} que todos los alumnos participantes cuenten previamente con su cuenta de AWS creada y activada bajo el \textbf{plan gratuito (Free Tier)}. Esto evitará retrasos operativos durante las sesiones y permitirá que los estudiantes se enfoquen directamente en la realización de los ejercicios prácticos.

\vspace{1.5cm}

\begin{flushright}
\textbf{Atentamente,}

\vspace{1.5cm}

\textbf{Emilio Izquierdo Montero} \\
Student Leader del AWS Student Group \\
Instituto Tecnológico de Matehuala
\end{flushright}
\end{document}
